\sezione{10}{Errori ricorrenti e controllo finale}{slide 145}

\begin{inbreve}
Gli errori negli schemi E-R si ripetono: scegliere male fra attributo ed entità o fra relationship ed entità, leggere le cardinalità al contrario, dimenticare o sbagliare gli identificatori. Questo capitolo li raccoglie e propone una checklist da usare prima di consegnare un esercizio.
\end{inbreve}

\subsection{I quattro errori più frequenti}

\begin{center}
\small
\renewcommand{\arraystretch}{1.4}\renewcommand\tabularxcolumn[1]{m{#1}}
\rowcolors{2}{chiaro}{white}
\begin{tabularx}{\linewidth}{>{\raggedright\headfont\bfseries\small}m{3.3cm} >{\raggedright\arraybackslash}X >{\raggedright\arraybackslash}m{4.4cm}}
\rowcolor{navy}\thd{Errore} & \thd{Domanda da farsi} & \thd{Dove ripassare}\\
Attributo o entità? & Il concetto ha proprietà proprie che ci interessano (la città con provincia e regione)? Se sì, entità collegata; altrimenti attributo. Mai entrambi. & Capitolo 5, «Relationship o attributo?»\\
Relationship o entità? & La stessa coppia di occorrenze può legarsi più volte (esami ripetuti, incarichi ripetuti)? Se sì, la relationship non basta. & Capitolo 3 (promozione a entità) e capitolo 6 (esame ripetuto)\\
Il «verso» delle cardinalità & La coppia è scritta accanto all'entità a cui si riferisce? $(1,1)$ accanto a Studente in Residenza: ogni studente risiede in una e una sola città. & Capitolo 5\\
Identificatori dimenticati o sbagliati & Ogni entità ne ha almeno uno? Quelli esterni passano solo per relationship $(1,1)$ (va bene anche $(0,1)$, ma lo trascuriamo)? & Capitolo 6\\
\end{tabularx}
\end{center}

\subsection{Checklist prima di consegnare}

\begin{passi}
  \item \hl{Concetti}: ogni entità ha occorrenze chiare? Ho distinto concetti diversi (opera, edizione, volume, copia)?
  \item \hl{Nomi}: singolari, espressivi, sostantivi per le relationship, unici nello schema?
  \item \hl{Relationship}: possono ripetersi le stesse coppie? Ho indicato i ruoli nelle ricorsive?
  \item \hl{Attributi}: ognuno sta sul costrutto giusto (il voto sull'esame, la data di restituzione sul prestito)?
  \item \hl{Cardinalità}: ho scritto minimo e massimo per ogni entità di ogni relationship, dal lato giusto?
  \item \hl{Identificatori}: almeno uno per entità, minimali, esterni solo con $(1,1)$?
  \item \hl{Generalizzazioni}: totali o parziali correttamente marcate? Nessuna entità figlia di due generalizzazioni?
  \item \hl{Verifica}: ho ripercorso lo schema con gli esempi delle specifiche?
  \item \hl{Documentazione}: dizionario dei dati e vincoli non esprimibili?
\end{passi}

\begin{mlbox}[Che cosa viene dopo]
Lo schema E-R è il punto di partenza della \key{progettazione logica}: entità e relationship diventano tabelle, identificatori diventano chiavi, identificatori esterni diventano chiavi che includono chiavi esterne. Le scelte fatte qui (in particolare gli identificatori) determinano direttamente lo schema relazionale.
\end{mlbox}

\begin{appunti}{Riepilogo del capitolo}
\pag{145} Attributo o entità; relationship o entità; verso delle cardinalità; identificatori.
\end{appunti}
